\documentclass{article}

\textwidth=6.5in
\textheight=8.5in
\voffset=-54pt
\oddsidemargin=0in
\evensidemargin=0in
\setlength{\parskip}{8pt}
\setlength{\parindent}{0pt}

\usepackage[normalem]{ulem}

\usepackage{amsmath, amssymb}
\usepackage{ifthen}

\usepackage[inline]{enumitem}
\setlist{topsep=0pt}

\usepackage{xcolor}
\usepackage[pdftex]{graphicx}

% change hyperref options
\PassOptionsToPackage{hyphens}{url}
\usepackage[bookmarks,colorlinks,linkcolor=black,citecolor=blue,pdfstartview=FitH,urlcolor=black]{hyperref}
\hypersetup{pdfpagemode=UseNone}

\usepackage{graphicx}
\usepackage{multicol}
\usepackage{framed}
\usepackage{array}
\usepackage{icomma}
\usepackage{microtype}
\usepackage{diffcoeff}[=v4]

\usepackage{soul}

\usepackage{pgfplots}
\pgfplotscreateplotcyclelist{mycolor}{%
  black, blue, red, green!70!black,magenta,
  purple, cyan, orange
}
\pgfplotsset{
  cycle list name=mycolor,
  axis lines=middle,
  every axis plot/.style={no marks, thick},
  every axis/.style={
    enlarge x limits=false,
    enlarge y limits=auto,
    xlabel style={at={(current axis.right of origin)},anchor=west},
    ylabel style={at={(current axis.above origin)},anchor=south},
    % extra x and y ticks for asymptotes
    extra x tick style={grid=major, grid style={dashed,black}},
    extra y tick style={grid=major, grid style={dashed,black}},
    /pgf/number format/1000 sep={},
  },
  tick label style = {font=\footnotesize, fill=white, inner sep=0pt},
  xticklabel shift=1pt,    
  scaled ticks=false,
  scaled y ticks=false,
  unbounded coords=jump,
  samples=101,
  minor tick length=0,
  scatter/classes={
    open={mark=*,draw=black,fill=white, mark size=2, thin},
    closed={mark=*,draw=black,fill=black, mark size=2, thin}
  },
  compat=1.18
}

\usepackage{tikz}
\usetikzlibrary{
  % enables the snake paths
  decorations.pathmorphing,%
  % enables bigger arrows
  decorations.markings,%
  decorations.pathreplacing,%
  decorations.text,%
  calc,%
  patterns,%
  shapes.geometric,%
  arrows,
  decorations.shapes,
  positioning,
  plotmarks,
}
\tikzset{>=latex} % sets default arrow type in tikz
\tikzset{font=\small}

\interfootnotelinepenalty=10000
\renewcommand{\thefootnote}{(\arabic{footnote})}

\usepackage{multirow, bigdelim}

\newcommand\iscurrentenumi[1]{%
  \ifthenelse{\equal{#1}{\theenumi}}{}{#1}}
\setenumerate[1]{ref=\#\arabic*}
\setenumerate[2]{ref=\protect\iscurrentenumi{\theenumi}(\alph*)}
\newcommand\enumiiprefix[2]{%
  \ifthenelse{{\equal{#1}{\theenumi}}\AND{\equal{#2}{(\alph{enumii})}}}{}{}%
  \ifthenelse{{\equal{#1}{\theenumi}}\AND{\NOT{\equal{#2}{(\alph{enumii})}}}}{#2}{}%
  \ifthenelse{\equal{#1}{\theenumi}}{}{#1#2}%
}
\setenumerate[3]{ref=\protect\enumiiprefix{\theenumi}{(\alph{enumii})}\roman*}

\definecolor{ideacolor}{rgb}{0.7,0,0}
\newcommand{\ideas}[1]{\color{ideacolor} #1 \color{black}}
\newcommand{\ideacolor}{maroon}

\newlength{\tipmargin}
\makeatletter
\newcommand{\version}[2]{\ifthenelse{\equal{#1}{1a}}{#2}{}}

\renewenvironment{framed}{
  \setlength{\tipmargin}{\textwidth-\linewidth}
  \advance\@totalleftmargin-\tipmargin
  \def\FrameCommand{\hspace{\tipmargin}\fbox}%
  \advance\linewidth\tipmargin
  \MakeFramed {\advance\hsize-\width \FrameRestore}%
}
{\endMakeFramed}

\def\pr@m@s{%
  \ifx'\@let@token
    \expandafter\pr@@@s
  \else
    \ifx^\@let@token
      \expandafter\expandafter\expandafter\pr@@@t
    \else
      \expandafter\expandafter\expandafter\pr@@@end
    \fi
  \fi}
\def\pr@@@end{\egroup
    \@ifnextchar({\mkern-2mu }{\@ifnextchar\{{\mkern-2mu }{\@ifnextchar[{\mkern-.5mu }{}}}}

\makeatother

\newlength{\highlightwidth}
\newcommand{\highlight}[1]{\setlength{\highlightwidth}{\linewidth}%
  \addtolength{\highlightwidth}{-55pt}%
  \hspace{24pt}\begin{minipage}[t]{\highlightwidth} \setlength{\parskip}{8pt} #1 \end{minipage}\hspace{24pt}}

\definecolor{purple}{rgb}{0.5,0,1.0}

\newcommand{\textbook}[0]{\href{http://isites.harvard.edu/fs/docs/icb.topic1511355.files/fullbook.pdf}{textbook}}

% change to \usepackage[sols]{handout} to see solutions
\usepackage[]{handout}
\renewcommand{\workshopnote}{CA Notes.}    

\begin{document}

\htitle{Worksheet 2 - The Chain Rule, Continued}
\begin{problemonly}
    \begin{framed}

You should be able to:

\begin{itemize}[topsep=0pt,partopsep=0pt,parsep=8pt,itemsep=0pt]

\item You should be able to apply the Chain Rule like a pro!

\item You should be able to apply the Chain Rule in more general contexts,
  such as in a word problem or when given a composition of functions
  described by graphs or numerical data.

\item You should understand how logarithms and the Chain Rule can be used
  to show that the Power Rule $\tfrac{d}{dx} \left( x^n \right) =
  nx^{n-1}$ holds for any real number $n$.  (In Math Ma, we only proved
  the Power Rule in the case where $n$ is an integer or $1/2$.)

\end{itemize} 

    \end{framed}
\end{problemonly}

\begin{enumerate}
\item 
  \note{This is practice applying the Chain Rule in context, as well as in
    interpreting derivatives.  This is going to be very important when we
    start related rates, so make sure students are comfortable doing
    something like this.}

  \begin{problem}[\stretch{2}]
    At a time we'll designate as $t = 0$, we start blowing up a spherical balloon. 
    We observe that, at time $t = 3$, the balloon has a
    radius of 5 cm, and the instantaneous rate of change of the balloon's
    \textbf{radius} with respect to time is 2 cm per second ($t$ is
    measured in seconds).  At this time, what is the instantaneous rate of
    change of the balloon's \textbf{volume} with respect to time?  (The
    volume of a sphere of radius $r$ is $\frac{4}{3} \pi r^3$.)  What are
    the units of your answer?    
  \end{problem}

  \begin{solution}
    Here are a couple of different ways to think about this problem.
    (They're really saying the same thing, just in different notation.)
    \begin{itemize}
    \item \textbf{Perspective 1.}  We're asked for the instantaneous rate
      of change of the balloon's volume with respect to time, so if we let
      $V$ be volume and $t$ be time, then we can express our goal as:

      \highlight{\textbf{Goal:} Find $\dfrac{dV}{dt}$ at time $t = 3$.}

      It's always a good idea to figure out what we know.  In this case,
      we're told a bunch of things about radius, so let's name that $r$.

      \highlight{\textbf{What we know:} When $t = 3$, $r = 5$ cm and
        $\dfrac{dr}{dt} = 2$ cm/s.}

      The Chain Rule tells us that the rate we want to know $\left(
        \dfrac{dV}{dt} \right)$ and the rate we do know $\left(
        \dfrac{dr}{dt} \right)$ are related by $\dfrac{dV}{dt} =
      \dfrac{dV}{dr} \dfrac{dr}{dt}$.  So, if we can find
      $\dfrac{dV}{dr}$, we'll be set.  We know that $V = \dfrac{4}{3} \pi
      r^3$, so $\dfrac{dV}{dr} = 4 \pi r^2$.  Now, we can put this
      together.  At time $t = 3$,
      \begin{align*}
        \frac{dV}{dt} &= \frac{dV}{dr} \frac{dr}{dt} \\
        &= (4 \pi r^2) (2) \\
        &= (4 \pi \cdot 5^2) \cdot 2 \\
        &= \boxed{200 \pi \text{ cm}^3/\text{s}}
      \end{align*}
      (We got the values for $\dfrac{dr}{dt}$ and $r$ from the ``What we
      know'' above.)

      The units must be cm$^3$/s because we are looking for a rate of
      change of volume (measured in cm$^3$) with respect to time (measured
      in s).

    \item \textbf{Perspective 2.}

      We know that the volume of a sphere of radius $r$ is $\dfrac{4}{3}
      \pi r^3$.  We can express this as $\displaystyle V(r) = \frac{4}{3}
      \pi r^3$.  However, in this problem, we're interested in how volume
      relates to \emph{time}.  We know that the radius is changing over
      time, so let $r(t) = $ radius at time $t$.  Then, the function that
      gives volume in terms of time is the composition $V(r(t))$, and we
      want to know the instantaneous rate of change change of this when $t
      = 3$.  That is:

      \highlight{\textbf{Goal:} Find $\dfrac{d}{dt} V(r(t))$ when $t =
        3$.  This is sometimes also written as $\left. \dfrac{d}{dt}
          V(r(t)) \right|_{t=3}$.}

      The Chain Rule tells us that:
      \begin{align*}
        \left. \frac{d}{dt} V(r(t)) \right|_{t=3} &= V'(r(3)) r'(3) \\
        \intertext{Since $V(r) = \dfrac{4}{3} \pi r^3$, $V'(r) = 4 \pi
          r^2$, so we can rewrite the previous line as:}
        &= 4 \pi \cdot r(3)^2 \cdot r'(3) \\
        \intertext{We are given that $r'(3) = 2$ cm/s and $r(3) = 5$ cm, 
          so we plug these in:}
        &= 4 \pi \cdot 5^2 \cdot 2 \\
        &= \boxed{200 \pi \text{ cm}^3/\text{s}}
      \end{align*}
    \end{itemize}
  \end{solution}



\item 
  \note{This is meant to show students why the Power Rule works for
    non-integer exponents.  It's also good practice with logs and
    exponents.}  

  In Math Ma, we learned the Power Rule, which said that $\frac{d}{dx} x^n
  = n x^{n-1}$ for any real number $n$.  However, we only \textit{proved}
  this when $n$ was either a positive integer or $\frac{1}{2}$.  Now that
  we know the Chain Rule, we can finally prove the Power Rule for all real
  numbers $n$.  In this problem, you'll use the Chain Rule to prove that
  the derivative of $x^{\pi}$ is $\pi x^{\pi - 1}$.  (The same idea works
  for any exponent, not just $\pi$.)

  \begin{enumerate}
  \item
    \begin{problem}[\vfill]
      Rewrite $f(x) = x^{\pi}$ as $f(x) = e^{\text{something}}$.  What
      is the ``something'' in the exponent?
    \end{problem}

    \begin{solution}
      Since $e^x$ and $\ln x$ are inverse functions, we can write
      \begin{align*}
        x^{\pi} &= e^{\ln (x^\pi)} \\
        &= e^{\pi \ln x}
      \end{align*}
    \end{solution}

  \item
    \begin{problem}[\vfill\newpage]
      Now, use the Chain Rule and what you know about the derivative of
      $e^x$ to show that $f'(x) = \pi x^{\pi - 1}$. 
    \end{problem}

    \begin{solution}
      By the Chain Rule,
      \begin{align*}
        \frac{d}{dx} \left(\fcolorbox{red}{white}{$\displaystyle e^ {
              \fcolorbox{blue}{white}{$\displaystyle \pi  \ln x$}} $}\right)
        = {} & e^ { \pi \ln x} \cdot \frac{d}{dx}
        \left(\fcolorbox{purple}{white}{$\displaystyle \pi$}\cdot \ln
          x\right) \\ 
        = {} & e^ { \pi \ln x} \cdot \frac{\pi}{x} \\
        = {} & x^\pi \cdot \frac{\pi}{x} \text{ since we saw in the previous
          part that $x^{\pi} = e^{\pi \ln x}$} \\
        = {} & \pi x^{\pi - 1}
      \end{align*}
      Alternatively, we could use logarithmic differentiation, which we
      will discuss in lesson 4. We start by taking the log of $f(x)$.
      \begin{align*}
        \ln f(x) &= \ln(x^\pi) \\
        \ln f(x) &= \pi \ln x 
      \end{align*}
      Then we can differentiate both sides and solve for $f'(x)$.
      \begin{align*}
        \diff*{\fcolorbox{red}{white}{$\ln \fcolorbox{blue}{white}
        {$f(x)$}$}}{x} &= \diff*{\pi\ln x}{x} \\
        \frac{1}{f(x)} \cdot f'(x) &= \pi \cdot \frac{1}{x} \\
        f'(x) &= \pi f(x) \cdot \frac{1}{x} \\
        &= \pi x^\pi \cdot \frac{1}{x} \\
        &= \pi x^{\pi-1}
      \end{align*}
    \end{solution}
  \end{enumerate}

\item 

  \note{Just practice differentiating more complicated functions.
    Students can sometimes get a little lost in these longer problems, so
    encourage them to write out extra steps if needed.  (They tend to want
    to just write down the entire derivative at once.)
    
    Matt note - I also like to encourage them to keep their focus on the structures (i.e. derivative of the inner function, etc.) as opposed to the specific numbers/values/variables. They can often get very caught up on what a solution looks like, even if they are confident in their knowledge of the structures at play.}

  Differentiate the following functions.  (Don't try to simplify your
  answers.) 

  \begin{enumerate}
  \item

    \begin{problem}[\vfill]
      $h(x) = e^{3x} + x^6 \cdot \ln (7x^2 + 3)$
    \end{problem}

    \begin{solution}
      \begin{align*}
        h'(x) = {} & 3 e^ { 3 x} + \frac{d}{dx} \left(\fcolorbox{green}{white}{$\displaystyle  x^ { 6} $}\cdot \fcolorbox{green}{white}{$\displaystyle  \ln\left(7 x^ { 2} + 3\right)$}\right) \\
        = {} & 3 e^ { 3 x} + x^ { 6} \cdot \frac{d}{dx} \left(\fcolorbox{red}{white}{$\displaystyle \ln\left(\fcolorbox{blue}{white}{$\displaystyle 7 x^ { 2} + 3$}\right)$}\right)+ 6 x^ { 5} \ln\left(7 x^ { 2} + 3\right) \\
        = {} & 3 e^ { 3 x} + x^ { 6} \cdot \frac{1}{7 x^ { 2} + 3}\cdot 14 x+ 6 x^ { 5} \ln\left(7 x^ { 2} + 3\right)
      \end{align*}
    \end{solution}

  \item
    \begin{problem}[\vfill]
      $j(u) = \ln (\ln u)$
    \end{problem}

    \begin{solution}
      $\displaystyle \frac{d}{du}
      \left(\fcolorbox{red}{white}{$\displaystyle
          \ln\left(\fcolorbox{blue}{white}{$\displaystyle \ln
              u$}\right)$}\right) = {} \frac{1}{\ln u}\cdot \frac{1}{u}$. 
    \end{solution}

  \end{enumerate}

\item 

  \note{Here, they'll need to remember that the derivative at a point can
    be visualized as slope of a tangent line.  This fact is always worth
    reinforcing!}
    
  \begin{problem}
    Here are graphs of two functions, $f(x)$ and $g(x)$.  If $h(x) =
    f(g(x))$, what is $h'(1)$?

    \begin{tabular}{ccc}
      $f(x)$ & & $g(x)$ \\ \\ 
      \includegraphics[width=1.75in]{Images/chain-rule-graphical-f.pdf} & 
      \hspace{0.25in} & 
      \includegraphics[width=1.75in]{Images/chain-rule-graphical-g.pdf} 
    \end{tabular}
  \end{problem}

  \begin{solution}
    By the Chain Rule, $h'(1) = f'(g(1)) g'(1)$.  From the graph of $g$,
    we see that $g(1) = -2$ and $g'(1) = -1$.  Therefore, $h'(1) =
    -f'(-2)$.  From the graph of $f$, we see that $f'(-2) = -3$, so
    $h'(1)$ is equal to \fbox{3}.
  \end{solution}

\item 
  \begin{problem}[\vfill\newpage]
    You are given the following information about three functions
    $f$, $g$, and $h$.
    \begin{align*}
      h(1) &= 2 & g(2) &= 3 & f'(3) &= 6 \\
      h'(1) &= 4 & g'(2) &= 5 
    \end{align*}
    If $r(x) = f(g(h(x)))$, do you have enough information to find
    $r'(1)$?  If so, compute it. If not, what additional information do
    you need?
  \end{problem}

  \begin{solution}
    First, let's find $r'(x)$:
    \begin{align*}
      r'(x) = {} & \frac{d}{dx}
      \left(\fcolorbox{red}{white}{$\displaystyle
          f\left(\fcolorbox{blue}{white}{$\displaystyle
              g\left(h(x)\right)$}\right)$}\right) \\
      = {} & f'\left(g\left(h(x)\right)\right)\cdot
      \frac{d}{dx} \left(\fcolorbox{red}{white}{$\displaystyle
          g\left(\fcolorbox{blue}{white}{$\displaystyle
              h(x)$}\right)$}\right) \\ 
      = {} & f'(g(h(x)))\cdot
      g'(h(x))\cdot h'(x)
    \end{align*}
    Plugging in $x = 1$, 
    \begin{align*}
      r'(1) &= f'(g(h(1))) g'(h(1)) h'(1) \\
      &= f'(g(2)) g'(2) \cdot 4 \\
      &= f'(3) \cdot 5 \cdot 4 \\
      &= 6 \cdot 5 \cdot 4 \\
      &= \boxed{120}
    \end{align*}
    (So, we did have enough information to find $r'(1)$.)
  \end{solution}

  \pspace{\vfill}

\item 
 % Gottlieb, 16.3, #7
  Here is the graph of a function $h(x)$.

  \probonly{
    \parbox[t]{2.5in}{\vspace{0pt}\includegraphics[width=2.5in]{Images/go16p09}}
    \hspace{0.1in}
    \begin{minipage}[t]{\linewidth-2.67in}
      \vspace{0pt}
    }

  \solonly{\includegraphics[width=2.5in]{Images/go16p09.pdf}}

  \begin{enumerate}
  \item
    \begin{problem}
      For which values of $x$ is $h'(x)$ negative?      
    \end{problem}

    \begin{solution}
      $h'(x)$ is negative whenever the graph of $h(x)$ is decreasing,
      which happens on the interval $(0, 3)$.
    \end{solution}
    
  \end{enumerate}
  \probonly{\end{minipage}}

  
  Define a new function $f(x)$ by $f(x) = [h(x)]^2$.
  \begin{enumerate}[resume]
    \item
      \begin{problem}[\vfill]
        Find all critical points of $f(x)$, and classify each as a
        local minimum, local maximum, or neither.  (\emph{Note:} This is
        asking about $f(x)$, not $h(x)$!)
      \end{problem}

    \begin{solution}
      We are looking for all $x$ such that $f'(x)$ is 0 or undefined.
      \version{Mb}{In addition, since the domain of $f$ is a closed
        interval, the endpoints $-3$ and $5$ will be critical points.}
      By the Chain Rule, $f'(x) = 2 h(x) h'(x)$.  This is never
      undefined, and it's equal to 0 when $h(x) = 0$ or $h'(x) = 0$.
      From the graph, $h(x) = 0$ when $x = -2$, $1$, or $4$, while
      $h'(x) = 0$ when $x = 0$ or $3$.  \version{1a}{So, the critical
        points of $f(x)$ are $\boxed{-2, 0, 1, 3, 4}$.}
      \version{Mb}{So, the critical points of $f(x)$ are $\boxed{-3, -2,
          0, 1, 3, 4, 5}$.}

      To classify these, let's use the First Derivative Test; that is,
      we'll make a sign chart for $f'$.  For example, let's figure out
      the sign of $f'(-2.5)$; this is equal to $2 h(-2.5) h'(-2.5)$;
      from the graph, we can see that $h(-2.5) < 0$ and $h'(-2.5) > 0$,
      so $f'(-2.5) < 0$.  The final sign chart looks like:
      \begin{center}
        \begin{tikzpicture}
          \begin{axis}[axis y line=none, x axis line style={-}, xtick={-3., -2., 0., 1., 3., 4., 5.}, xticklabels={$-3$, $-2$, $0$, $1$, $3$, $4$, $5$}, ymin=-2, ymax=2, height=1.25in, width=4in]
            \addplot+[domain=-7.:5.]{0};
            \draw (axis cs: -7., -1.5) node[right] {sign of $f'$};
            \draw (axis cs: -7., 1) node[right] {$f$};
            \draw (axis cs: -2.5, -1.5) node{$-$};
            \draw[->, xshift=2pt] (axis cs: -3, 1.5) -- (axis cs: -2.26667, 0.5);
            \draw (axis cs: -1., -1.5) node{$+$};
            \draw[->, xshift=2pt] (axis cs: -1.73333, 0.5) -- (axis cs: -0.266667, 1.5);
            \draw (axis cs: 0.5, -1.5) node{$-$};
            \draw[->, xshift=2pt] (axis cs: 0.266667, 1.5) -- (axis cs: 0.733333, 0.5);
            \draw (axis cs: 2., -1.5) node{$+$};
            \draw[->, xshift=2pt] (axis cs: 1.26667, 0.5) -- (axis cs: 2.73333, 1.5);
            \draw (axis cs: 3.5, -1.5) node{$-$};
            \draw[->, xshift=2pt] (axis cs: 3.26667, 1.5) -- (axis cs: 3.73333, 0.5);
            \draw (axis cs: 4.5, -1.5) node{$+$};
            \draw[->, xshift=2pt] (axis cs: 4.26667, 0.5) -- (axis cs: 5, 1.5);
            \draw[->, xshift=2pt] (axis cs: -2.26667, 0.5) -- (axis cs: -1.73333, 0.5);
            \draw[->, xshift=2pt] (axis cs: -0.266667, 1.5) -- (axis cs: 0.266667, 1.5);
            \draw[->, xshift=2pt] (axis cs: 0.733333, 0.5) -- (axis cs: 1.26667, 0.5);
            \draw[->, xshift=2pt] (axis cs: 2.73333, 1.5) -- (axis cs: 3.26667, 1.5);
            \draw[->, xshift=2pt] (axis cs: 3.73333, 0.5) -- (axis cs: 4.26667, 0.5);
          \end{axis}
        \end{tikzpicture}
      \end{center}
      From the sign chart, we see that $f$ has \fbox{local maxima at $x
        = 0, 3$} \\ and \fbox{local minima at $x = -2, 1, 4$}. 
    \end{solution}

    \item
    \begin{problem}[\vfill]
      Does $f$ achieve an absolute maximum on $[-3, 5]$?  If so, where?
    \end{problem}

    \begin{solution}
      Since $f$ is continuous and $[-3, 5]$ is a closed interval, $f$
      must have an absolute maximum on $[-3, 5]$ by the Extreme Value
      Theorem.  From our sign chart for $f'$, we can see that the
      absolute maximum must occur at $x = -3$, $x = 0$, $x = 3$, or $x =
      5$; to decide which one is correct, we can simply test the value
      of $f$ at all these points:
      \begin{align*}
        f(-3) &= h(-3)^2 = (-4)^2 = 16 \\
        f(0) &= h(0)^2 = 3^2 = 9 \\
        f(3) &= h(3)^2 = (-3)^2 = 9 \\
        f(5) &= h(5)^2 = 2^2 = 4
      \end{align*}
      So, $f$ has its absolute maximum at $\boxed{x = -3}$. 
    \end{solution}

  \item
    \begin{problem}[\vfill]
      Does $f$ achieve an absolute minimum on $[-3, 5]$?  If so, where?
    \end{problem}

    \begin{solution}
      Since $f$ is continuous and $[-3, 5]$ is a closed interval, $f$
      must have an absolute minimum on $[-3, 5]$ by the Extreme Value
      Theorem.  From our sign chart for $f'$, we can see that the
      absolute minimum must occur at $x = -2$, $x = 1$, or $x = 4$.
      Let's test the value of $f$ at these points:
      \begin{align*}
        f(-2) &= h(-2)^2 = 0^2 = 0 \\
        f(1) &= h(1)^2 = 0^2 = 0 \\
        f(4) &= h(4)^2 = 0^2 = 0
      \end{align*}
      So, the absolute minimum of $f$ on $[-3, 5]$ occurs at 3 points,
      $\boxed{x = -2, 1, 4}$. 
    \end{solution}

  \end{enumerate}    

\end{enumerate}

\end{document}